\section{What remains open}
\label{sec:open}

\paragraph{Unlocked access to shared state.}
The dissertation opens with Dekker's algorithm, made correct on this
machine by a fence after each store. Each thread reads a flag that the
other writes, with no lock. The calculus cannot express that: a thread may
touch a cell only if it owns it or holds the lock. The dissertation's basic ``daring'' rules were meant for this. They would let a
thread load from or store to the invariant's cells without the lock,
provided that the single command, specified with $J$ in hand, gives $J$
back. The dissertation doubted their soundness and left them as future
work. For the repaired calculus the doubt was justified. Adding the load
rule lets one derive a false specification
(\rocq{daring_load_extension_breaks_adequacy}), and each rule has an
instance whose premise is valid and whose conclusion is invalid
(\rocq{daring_load_honest_counterexample},
\rocq{daring_store_honest_counterexample}). The store rule fails even
while the lock is free. A sound rule for unlocked access, and with it a
proof of Dekker's algorithm, is still missing.

\paragraph{Realistic locks.}
The one lock is global and blocks every other processor. Programmer-level
locks can be encoded as flag cells under the invariant, as the container
does, but each release must itself take the global lock. A spinlock
released by a bare store is another instance of unlocked access.

\paragraph{A strong temporal connective.}
The frame property holds for a single split that is both spatial and
temporal (\rocq{owned_safe_common_frame}). Making that a connective of the
assertion language, in place of the conjunction of two, would restore
\rn{frame-stm}. Whether its other laws would then hold is not known.

\paragraph{Assertions.}
The proof system for assertions is sound and was sufficient for every
derivation attempted. Nothing is known about what it cannot prove, and
nothing is automated.

\paragraph{The rest of the dissertation.}
Allocation, recursive predicates, permissions and the other topics of
Chapter~5 are untouched, as is the sequential logic of Chapter~3.
